%% kapa: exam=1 round=30 year=2019 period=1 form=A subject=real_estate subject_ko=부동산학원론 seq=26 number=106
%% source: Q-Net 감정평가사 기출문제 (cst00309) artlSeq=5205205 file=extracted/1008/5205205/2019년 제30회 감정평가사 1차시험 문제지/2019년 제30회 감정평가사 1차시험 문제지/2019년도 제30회 감정평가사 1차 1교시 A형.hwp
%% conversion: kordoc 4.24.51 -> tools/gen1.py (auto); review=proofread-sonnet; flags=none
\begin{question}{106}
주거분리와 여과과정에 관한 설명으로 옳지 않은 것은?
\begin{choices}
  \item 저가주택이 수선되거나 재개발되어 상위계층의 사용으로 전환되는 것을 상향여과라 한다.
  \item 민간주택시장에서 저가주택이 발생하는 것은 시장이 하향여과작용을 통해 자원할당기능을 원활하게 수행하고 있기 때문이다.
  \item 주거입지는 침입과 천이현상으로 인해 변화할 수 있다.
  \item 주거분리는 도시 전체에서 뿐만 아니라 지리적으로 인접한 근린지역에서도 발생할 수 있다.
  \item 하향여과는 고소득층 주거지역에서 주택의 개량을 통한 가치상승분이 주택개량비용보다 큰 경우에 발생한다.
\end{choices}
\end{question}
\answer{5}
